%%%%%%%%%%%%%%%%%%%%%%%%%%%%%%%%%%%%%%%%%%%%%%%%%%%%%%%%%%%%%%%%%%%%%%%%%%%%%%%%
%  Policy Evaluation Lab -- Lab 4: a case study, the PNRR Parco agrisolare
%  Friday 6 November 2026 (hidden until 17:00). In the original table the region names are misaligned with the rows (wrapped labels); rows here follow the standard order of the regions, which the original label list also follows.
%  Source: MEF written exam, 31 March 2025, code STAT, envelope (busta) 1
%  (not drawn), question 2; published by the MEF on 4 June 2025 (ZIP of the
%  written-test papers on the competition page). Case text and questions
%  translated; Table 1 re-typed and checked: every column adds up to the
%  published total. Data constructed for the exercise (note kept).
%%%%%%%%%%%%%%%%%%%%%%%%%%%%%%%%%%%%%%%%%%%%%%%%%%%%%%%%%%%%%%%%%%%%%%%%%%%%%%%%

\documentclass[aspectratio=169,xcolor=table]{beamer}
\input{preamble-pel}
\newcolumntype{L}[1]{>{\raggedright\arraybackslash}p{#1}}

\title[Policy Evaluation Lab]{Policy Evaluation Lab}
\subtitle[Lab 4]{Lab 4: A case study, the PNRR \textit{Parco agrisolare}}
\date[\pelYear]{November, 2026}

\begin{document}

\frame{\titlepage}

%%%%%%%%%%%%%%%%%%%%%%%%%%%%%%%%%%%%%%%%%%%%%%%%%%%%%%%%%%%%%%%%%%%%%%%%%%%%%%%%
%  Italy's PNRR slides: every statement and number taken from the Council
%  document ST 12357/26 ADD 1 (20 August 2026), Annex to the Council
%  Implementing Decision amending the implementing decision of 13 July 2021
%  on the approval of the assessment of Italy's recovery and resilience plan
%  (Italian version; file Annex_CONSIL_ST_12357_2026_ADD_1_IT_TXT.pdf in the
%  repository). Pages: total cost p. 438; instalments, Section 2; Parco
%  agrisolare pp. 172, 177-179, 183, 195 (document page numbers).
%  The EU-level slides from next-generation-eu.europa.eu are kept in
%  lab-04-before-italy.tex (scratchpad) / git history.
%%%%%%%%%%%%%%%%%%%%%%%%%%%%%%%%%%%%%%%%%%%%%%%%%%%%%%%%%%%%%%%%%%%%%%%%%%%%%%%%
\section{Italy's recovery plan (PNRR)}

\begin{frame}{The PNRR in one document}
\bi
  \item The Council approved the assessment of Italy's \textbf{Piano Nazionale di Ripresa e Resilienza} (PNRR) with an Implementing Decision of \textbf{13 July 2021}
  \item That decision has been amended: this Annex, dated \textbf{20 August 2026}, belongs to an amending decision
  \item The Annex lists, for every reform and investment, its description, its \textbf{milestones and targets}, and the instalments to which they are linked
\ei
\vspace{0.15cm}
\begin{block}{Estimated total cost of the plan}
\centering \textbf{€194,435,381,164}\\[2pt]
\small of which the REPowerEU chapter: €7,212,200,000
\end{block}
\pelTableNote{Source: Council of the EU, ST 12357/26 ADD 1, 20 August 2026 (Annex on Italy's recovery and resilience plan).}
\end{frame}

\begin{frame}{Missions and components}
\begin{columns}[T]
\begin{column}{0.5\textwidth}
\scriptsize\setlength{\itemsep}{0pt}
\textbf{Mission 1}
\begin{itemize}\setlength{\itemsep}{0pt}\setlength{\parskip}{0pt}
  \item C1: digitalisation of the public administration, justice, public administration
  \item C2: digitalisation, innovation and competitiveness of the production system (axis 1)
  \item C3: tourism and culture 4.0
\end{itemize}
\textbf{Mission 2}
\begin{itemize}\setlength{\itemsep}{0pt}\setlength{\parskip}{0pt}
  \item C1: \textbf{circular economy and sustainable agriculture}
  \item C2: energy transition and sustainable mobility
  \item C3: energy efficiency and renovation of buildings
  \item C4: protection of land and water resources
\end{itemize}
\textbf{Mission 3}
\begin{itemize}\setlength{\itemsep}{0pt}\setlength{\parskip}{0pt}
  \item C1: infrastructure for sustainable mobility
\end{itemize}
\end{column}
\begin{column}{0.48\textwidth}
\scriptsize\setlength{\itemsep}{0pt}
\textbf{Mission 4}
\begin{itemize}\setlength{\itemsep}{0pt}\setlength{\parskip}{0pt}
  \item C1: education services, from nurseries to university
  \item C2: from research to business
\end{itemize}
\textbf{Mission 5}
\begin{itemize}\setlength{\itemsep}{0pt}\setlength{\parskip}{0pt}
  \item C1: labour policies
  \item C2: social infrastructure, families, communities and third sector
  \item C3: special interventions for territorial cohesion
\end{itemize}
\textbf{Mission 6}
\begin{itemize}\setlength{\itemsep}{0pt}\setlength{\parskip}{0pt}
  \item C1: proximity networks, facilities and telemedicine for territorial health care
  \item C2: innovation, research and digitalisation of the national health service
\end{itemize}
\textbf{Mission 7}: REPowerEU
\end{column}
\end{columns}
\pelTableNote{Source: Council of the EU, ST 12357/26 ADD 1, 20 August 2026 (Annex on Italy's recovery and resilience plan).}
\end{frame}

\begin{frame}{Paid by results: milestones, targets, instalments}
\begin{columns}[T]
\begin{column}{0.5\textwidth}
\footnotesize
\bi
  \item Italy receives the money in \textbf{instalments}: 10 of non-repayable support (grants) and 10 of loans
  \item Each instalment is linked to a list of:
    \bi
      \item \textbf{milestones} (\textit{traguardi}): qualitative steps, e.g.\ a law entering into force
      \item \textbf{targets} (\textit{obiettivi}): quantitative indicators, with a unit, a baseline, a value to reach and a quarter
    \ei
  \item Example, 1st grant instalment: entry into force of the primary legislation on the governance of the PNRR
\ei
\end{column}
\begin{column}{0.46\textwidth}
\centering\scriptsize
\begin{tabular}{lrr}
\toprule
\textbf{Instalment} & \textbf{Grants (€ m)} & \textbf{Loans (€ m)} \\
\midrule
1st  & 11,494.3 & 12,643.7 \\
2nd  & 11,494.3 & 12,643.7 \\
3rd  & 11,494.3 & 9,825.3 \\
4th  & 2,315.6  & 16,611.5 \\
5th  & 3,621.5  & 9,030.6 \\
6th  & 2,053.6  & 7,908.5 \\
7th  & 5,294.6  & 15,716.0 \\
8th  & 3,541.6  & 11,184.6 \\
9th  & 7,670.1  & 7,118.5 \\
10th & 12,799.9 & 19,919.6 \\
\bottomrule
\end{tabular}
\pelTableNote{Amounts rounded to €0.1 million from the Annex, Section 2.}
\end{column}
\end{columns}
\pelTableNote{Source: Council of the EU, ST 12357/26 ADD 1, 20 August 2026 (Annex on Italy's recovery and resilience plan).}
\end{frame}

\begin{frame}{Today's measure in the plan: Parco agrisolare}
\small
\textbf{Mission 2, Component 1, Investment 2.2 -- Parco agrisolare}: to make the agricultural sector more sustainable and competitive, by supporting investments in agricultural, livestock and agro-industrial buildings for the installation of \textbf{solar panels} and complementary interventions
\vspace{0.1cm}

\centering\scriptsize
\begin{tabular}{llL{6.2cm}r}
\toprule
\textbf{Code} & \textbf{Type} & \textbf{What must be reached} & \textbf{By} \\
\midrule
M2C1-4 & target & beneficiary projects identified, worth at least 30\% of the resources assigned to the investment & Q4 2022 \\
M2C1-5 & target & at least 32\% & Q4 2023 \\
M2C1-6 & target & at least 63.5\% & Q2 2024 \\
M2C1-6 bis & target & 100\% of the further supplementary resources & Q4 2024 \\
M2C1-9 & target & \textbf{at least 1,470,000 kW (1,470 MW)} of solar capacity installed & \textbf{Q2 2026} \\
\bottomrule
\end{tabular}
\pelTableNote{The exam text (2025) gives a target of 1,383 MW: the current Annex says 1,470 MW.}
\pelTableNote{Source: Council of the EU, ST 12357/26 ADD 1, 20 August 2026 (Annex on Italy's recovery and resilience plan).}
\end{frame}

\begin{frame}{A second instrument: the Parco agrisolare facility}
\small
\textbf{Mission 2, Component 1, Investment 4 -- Dispositivo per il parco agrisolare}
\bi
  \item A grant scheme to encourage private investment and improve access to finance in the agricultural and agro-industrial sector
  \item Managed by the \textbf{GSE} (Gestore dei Servizi Energetici S.p.A.); it supports the same kind of investments: solar panels on agricultural, livestock and agro-industrial buildings
  \item Italy transfers \textbf{€789 million} to the GSE for the scheme; at least €633 million must contribute to climate objectives
  \item Milestones: implementation agreement in force (M2C1-26, Q4 2025); legal agreements with final beneficiaries for 100\% of the scheme (M2C1-27, Q2 2026)
\ei
\pelTableNote{Source: Council of the EU, ST 12357/26 ADD 1, 20 August 2026 (Annex on Italy's recovery and resilience plan).}
\end{frame}

\section{Today's case}

\begin{frame}{So, today we are going to evaluate a case study}
\be
  \item We look at a real case study
  \item Form groups of \textbf{2--3 people} at most
  \item Study what the question is asking, and look at the data
  \item Formulate some opinions, and the approach you would use
  \item Present your results
\ee
\end{frame}

\begin{frame}{Today's lab: a case study}
\small
Italy is among the countries with the highest direct energy consumption in food production in the European Union. A significant part of the investments of the National Recovery and Resilience Plan (PNRR) is meant to support \textbf{solar photovoltaic plants} in the agricultural and agro-industrial sector, promoting sustainable energy solutions for firms.

\vspace{0.15cm}
The \textbf{``Parco agrisolare''} investment, for which a total of \textbf{€2.35 billion} of PNRR funds has been allocated, aims to reduce energy consumption by upgrading production buildings and installing photovoltaic panels on the roofs of agricultural buildings (e.g.\ stables), with the target of reaching an installed capacity of at least \textbf{1,383 MW by 2026}.
\end{frame}

\begin{frame}{Table 1: Parco agrisolare by region}
\centering\footnotesize
\resizebox{!}{0.63\textheight}{%
\begin{tabular}[b]{l rrrrr}
\toprule
\textbf{Region} & \textbf{Projects} & \textbf{Total cost} & \textbf{PNRR funding} & \textbf{Planned capacity} & \textbf{Installed capacity} \\
 & \textbf{activated} & \textbf{of projects} & \textbf{assigned} & \textbf{(MW, at the end} & \textbf{(MW, completed} \\
 & \textbf{at 31/10/2024} & \textbf{(€ million)} & \textbf{(€ million)} & \textbf{of the works)} & \textbf{at 31/10/2024)} \\
\midrule
Liguria & 87 & 9.7 & 5.6 & 4.9 & 1.0 \\
Lombardia & 2,038 & 500.2 & 283.6 & 247.5 & 44.9 \\
Piemonte & 1,348 & 236.1 & 133.3 & 119.1 & 21.8 \\
Valle d'Aosta & 29 & 2.3 & 1.2 & 0.7 & 0.1 \\
Emilia-Romagna & 1,520 & 340.1 & 181.3 & 159.9 & 25.8 \\
Friuli-Venezia Giulia & 427 & 81.3 & 46.1 & 41.2 & 5.4 \\
Trentino-Alto Adige & 856 & 96.2 & 55.3 & 37.8 & 7.4 \\
Veneto & 1,641 & 295.8 & 175.0 & 152.6 & 23.1 \\
Lazio & 524 & 76.2 & 45.1 & 38.3 & 2.6 \\
Marche & 327 & 49.9 & 26.3 & 26.1 & 3.4 \\
Toscana & 844 & 118.1 & 66.8 & 51.7 & 5.7 \\
Umbria & 282 & 41.8 & 21.8 & 22.2 & 4.2 \\
Abruzzo & 348 & 51.9 & 31.5 & 26.3 & 4.4 \\
Basilicata & 237 & 36.7 & 23.4 & 18.6 & 1.5 \\
Calabria & 553 & 64.7 & 42.1 & 32.3 & 2.8 \\
Campania & 815 & 126.4 & 78.0 & 60.8 & 8.5 \\
Molise & 196 & 31.5 & 20.0 & 15.0 & 1.7 \\
Puglia & 1,334 & 171.0 & 111.1 & 82.5 & 10.2 \\
Sardegna & 544 & 70.9 & 41.0 & 28.6 & 1.8 \\
Sicilia & 1,135 & 151.1 & 94.0 & 70.3 & 6.4 \\
\midrule
\textbf{Italy} & \textbf{15,085} & \textbf{2,551.9} & \textbf{1,482.5} & \textbf{1,236.4} & \textbf{182.7} \\
\bottomrule
\end{tabular}}
\pelTableNote{Solar plants on agricultural and livestock buildings. Source: MEF written exam, 31 March 2025 (envelope 1, not drawn); region names re-aligned with the rows, which are misprinted in the original. The data were constructed for the exercise and do not necessarily correspond to real values.}
\end{frame}

\begin{frame}{The questions}
\begin{block}{}
Based on the information provided, analyse the main evidence describing the implementation process of the investment, indicating also which indicators it would be useful to compute, and how.
\end{block}
\vspace{0.2cm}
\begin{block}{}
In addition, illustrate which strategy could be adopted, and which other data would be needed, to evaluate the effectiveness of the measure.
\end{block}
\end{frame}

\end{document}
